\section{상세 데이터 생성 방법}
\label{app:data_generation}

본 연구는 웹 스크린샷의 정보를 서로 다른 포괄성 수준으로 언어화하기
위해 2단계 합성 파이프라인을 사용한다. Stage 1에서는 스크린샷과
구조화된 웹 정보를 바탕으로 화면에서 검증 가능한 evidence를 추출하고,
Stage 2에서는 동일한 evidence를 서로 다른 information budget으로
언어화하여 Info-S, Info-M, Info-L을 생성한다. 세 조건은 동일한
스크린샷과 Stage 1 출력을 공유한다. 따라서 조건 간 차이는 사실 추출
과정이 아니라, 추출된 정보 중 어느 범위를 학습 target에 포함하는지에서
발생한다.

\subsection{입력과 Stage 1 Evidence 추출}

각 예제의 입력은 현재 viewport의 스크린샷, DOM tree, accessibility
tree로 구성된다. 스크린샷은 실제 시각 상태를 판단하는 주된 근거이며,
DOM과 accessibility tree는 텍스트, 역할, 상태, 클릭 가능 여부와
bounding box를 확인하는 보조 근거다. DOM에 존재하더라도 화면 밖에
있거나 다른 레이어에 가려진 요소는 현재 가시 evidence로 간주하지
않는다.

Stage 1은 다음 항목을 구조화한다.
\begin{itemize}
    \item \textbf{Page type and purpose}: 검색, 로그인, 상품 탐색,
          예약 또는 설정 등 현재 인터페이스의 유형과 목적,
    \item \textbf{Visible facts}: 제목, 본문, 메뉴, 라벨, 가격, 날짜,
          수치, 경고, 선택값과 입력값,
    \item \textbf{Interface state}: 선택, 비활성화, 체크, 펼침, 입력,
          팝업과 같은 현재 상태,
    \item \textbf{Affordances}: 클릭, 입력, 선택 및 제출이 가능한
          요소의 이름, 행동 유형, 좌표와 예상 결과.
\end{itemize}
행동의 예상 결과는 아직 관찰되지 않은 예측이므로 현재 화면의 사실과
분리하여 저장하며, 확인할 수 없는 목적지나 후속 내용을 생성하지 않는다.

\subsection{가시성 및 팝업 처리}

구조화 트리는 화면 밖 요소, 숨겨진 요소와 팝업에 가려진 배경 요소를
포함할 수 있다. 이를 줄이기 위해 모델은 스크린샷을 보고 DOM 텍스트가
현재 보이는지 판단한다. 배경을 가리고 상호작용을 차단하는 모달이
표시되면 모달의 내용과 행동을 우선하며, 가려진 배경 요소를 현재 사용
가능한 affordance로 기술하지 않는다. 반면 지도나 검색 레이어처럼
레이어 자체가 화면의 주된 기능이면 이를 닫아야 하는 방해 요소로
취급하지 않는다. 배너, 토스트와 작은 메뉴는 현재 상태나 가능한 행동을
바꿀 때만 포함한다.

\subsection{좌표 추출과 정규화}

합성 모델이 픽셀만 보고 좌표를 추정하지 않도록 DOM 또는 accessibility
tree의 bounding box를 사용한다. 이미지의 너비와 높이가 $W,H$이고
요소의 bounding box가 $(x,y,w,h)$일 때 중심점을 다음과 같이
$0$--$999$ 좌표계로 변환한다.
\begin{equation}
\hat{x}=\operatorname{clip}\!\left(
\operatorname{round}\left(\frac{x+w/2}{W}\times999\right),0,999\right),
\end{equation}
\begin{equation}
\hat{y}=\operatorname{clip}\!\left(
\operatorname{round}\left(\frac{y+h/2}{H}\times999\right),0,999\right).
\end{equation}
Stage 2는 Stage 1에서 제공한 좌표를 그대로 인용하며 새 좌표를 추정,
반올림 또는 수정하지 않는다.

\subsection{Stage 2: 포괄성 수준별 Verbalization}

Stage 2는 스크린샷과 공통 Stage 1 evidence를 하나의 plain-text
description으로 변환한다. 세 조건 모두 동일한 형식을 사용하되,
target에 담는 평균 evidence coverage를 달리한다. 각 sample에서
선택된 evidence 집합의 엄격한 포함관계는 강제하지 않지만, 페이지
목적과 primary task flow는 모든 조건에서 우선한다.

\paragraph{Info-S (Concise).}
3--5문장과 약 90--130단어로 페이지 유형, 주된 목적, 가장 중요한
가시 내용 또는 상태와 최대 두 개의 interaction goal을 기술한다.
화면 전체를 열거하지 않고 화면을 대표하는 salient evidence를
선택한다. 행동 요소가 언급되면 해당 Stage 1 좌표를 함께 제공한다.

\paragraph{Info-M (Balanced).}
6--9문장과 약 170--230단어로 페이지 목적과 현재 상태, 각 주요
section의 대표 사실과 primary task를 기술한다. 약 3--4개의
interaction goal과 최대 6개의 actionable element를 포함하며, 중앙
입력창, 제출 버튼과 주요 CTA를 utility navigation보다 우선한다.
표와 반복 목록은 header와 대표 값을 기술하고, 핵심 행동 1--2개의
예상 결과만 추측형으로 표현한다.

\paragraph{Info-L (Comprehensive).}
Stage 1의 visible facts와 affordances를 가능한 한 빠짐없이 하나의
자기완결적 장문으로 변환한다. 세부 텍스트, 수치, 상태, section,
요소 기능, 좌표, 가능한 행동과 예상되는 상태 변화를 함께 기술하고,
표와 목록도 개별 항목과 값을 보존한다. 관찰된 사실은 사실형으로,
행동 결과는 ``may'', ``is expected to''와 같은 불확실성 표현으로
구분한다. 공간, 기능, 사용자 목표 또는 콘텐츠 순서 등 전개 방식은
달라질 수 있지만 포함해야 하는 evidence의 범위는 유지한다.

\begin{table}[t]
\centering
\resizebox{\columnwidth}{!}{%
\begin{tabular}{lccc}
\toprule
 & Info-S & Info-M & Info-L \\
\midrule
Visible evidence & Salient & Representative & Near-complete \\
Interaction goals & $\leq2$ & 3--4 & Near-complete \\
Actionable elements & Selected & $\leq6$ & Near-complete \\
Expected outcomes & Minimal & Key actions & Most actions \\
Target length & 90--130 & 170--230 & Unbounded \\
\bottomrule
\end{tabular}
}
\caption{세 webpage verbalization 조건의 목표 information budget.
길이 단위는 영문 단어 수이며, Info-L은 완전성을 우선하므로 고정된
길이 상한을 두지 않는다.}
\label{tab:appendix_info_levels}
\end{table}

이 비교에서 길이는 독립적으로 조작된 변수가 아니다. 더 많은
evidence를 언어화할수록 target token 수도 함께 증가한다. 따라서
Info-S--Info-M--Info-L은 길이만 다른 caption이 아니라 서로 다른
information coverage를 갖는 실용적인 supervision recipe다.

\subsection{자동 검증과 Rejection}

Stage 1과 Stage 2 출력에는 다음 검사를 적용한다.
\begin{itemize}
    \item 출력 좌표가 해당 Stage 1 좌표 집합에 포함되는지 확인한다.
    \item 행동 가능하다고 기술된 요소에 정확한 좌표가 있는지 확인한다.
    \item Stage 1에 없는 텍스트, 값 또는 요소가 생성되지 않았는지
          확인한다.
    \item 별개의 라벨이나 메뉴가 새로운 복합 개념으로 병합되지
          않았는지 확인한다.
    \item 예상 결과가 관찰된 사실처럼 단정되지 않았는지 확인한다.
    \item Info-S와 Info-M의 단어, 문장 및 요소 budget을 확인한다.
    \item Info-L의 문장 및 evidence 반복 퇴화를 확인한다.
\end{itemize}
검증에 실패한 출력은 제한된 횟수만큼 재생성하며, 이후에도 좌표 발명,
형식 위반 또는 반복이 남은 예제는 최종 집합에서 제거한다.

\subsection{통제된 비교 조건}

Description supervision의 추가 효과와 information budget의 효과를
분리하기 위해 Grounding-only, Info-S, Info-M, Info-L의 네 조건을
구성한다. Description 조건에는 모두 공통 grounding 데이터를 혼합하고,
모두 생성에 성공한 동일한 screenshot 교집합을 사용한다. 이미지 수,
학습 update, backbone과 후속 instruction tuning, action warm-up 및
trajectory training은 동일하게 유지한다. Grounding-only와 각 Info
조건의 차이는 verbal supervision의 추가 효과를, Info-S--Info-M--Info-L
차이는 증가하는 평균 information budget의 효과를 나타낸다.
